\documentclass[10pt]{beamer}

\usetheme{Madrid}

\usepackage{amsmath,amssymb,bm}
\usepackage{physics}
\usepackage{graphicx}

\setbeamersize{text margin left=8mm,text margin right=8mm}
\setlength{\abovedisplayskip}{4pt}
\setlength{\belowdisplayskip}{4pt}
\setlength{\abovedisplayshortskip}{3pt}
\setlength{\belowdisplayshortskip}{3pt}

\title[Mass Scaling of HO Matrix Elements]{Mass Scaling of Harmonic Oscillator Matrix Elements in Nuclear Physics}
\author{Morten Hjorth-Jensen}
\date{Notes from March 1996}

\begin{document}

%------------------------------------------------
\begin{frame}
\titlepage
\end{frame}

%------------------------------------------------
\begin{frame}[t]{Motivation}
\small
In shell-model calculations, effective two-body matrix elements are often
constructed for a specific mass region.

\medskip

\begin{block}{Problem}
Given matrix elements fitted or computed for a nucleus with mass number
\(A_0=100\), how can they be approximately scaled to a nucleus with
\(A=132\)?
\end{block}

\medskip

The answer is connected to:

\begin{itemize}
    \item harmonic oscillator basis functions,
    \item oscillator frequency \(\hbar\omega\),
    \item oscillator length parameter \(b\),
    \item phenomenological shell-model scaling laws.
\end{itemize}
\end{frame}

%------------------------------------------------
\begin{frame}[t]{Harmonic Oscillator Basis}
\small
Single-particle shell-model states are commonly expanded in a harmonic
oscillator basis. The harmonic oscillator Hamiltonian is
\[
H=\frac{p^2}{2m}+\frac{1}{2}m\omega^2 r^2 .
\]

The corresponding oscillator length is
\[
b=\sqrt{\frac{\hbar}{m\omega}} .
\]

The radial wave functions therefore depend explicitly on the oscillator
frequency. Changing the mass number \(A\) changes the optimal value of
\(\hbar\omega\) and thus the spatial extent of the basis.
\end{frame}

%------------------------------------------------
\begin{frame}[t]{Empirical Oscillator Frequency}
\small
A widely used empirical formula is
\[
\hbar\omega \approx 41 A^{-1/3}\ \mathrm{MeV} .
\]

A more refined formula is
\[
\hbar\omega \approx 45A^{-1/3}-25A^{-2/3}\ \mathrm{MeV} .
\]

For the nuclei considered here:
\[
\hbar\omega(100)\approx 8.53\ \mathrm{MeV}, \qquad
\hbar\omega(132)\approx 7.87\ \mathrm{MeV} .
\]

Hence
\[
\frac{\hbar\omega(132)}{\hbar\omega(100)}\approx 0.923 .
\]
\end{frame}

%------------------------------------------------
\begin{frame}[t]{Oscillator Length Scaling}
\small
Since
\[
b=\sqrt{\frac{\hbar}{m\omega}},
\]
we obtain
\[
b\propto \omega^{-1/2} .
\]

Using the empirical scaling \(\omega\propto A^{-1/3}\), we find
\[
b\propto A^{1/6} .
\]

Therefore
\[
\frac{b_{132}}{b_{100}}
=\left(\frac{132}{100}\right)^{1/6}
\approx 1.047 .
\]
\end{frame}

%------------------------------------------------
\begin{frame}[t]{Interpretation of Oscillator-Length Scaling}
\small
The scaling
\[
\frac{b_{132}}{b_{100}}\approx 1.047
\]
means that the oscillator basis for \(A=132\) is about \(4.7\%\) more
spatially extended than that for \(A=100\).

\medskip

Interpretation:
\begin{itemize}
    \item The basis for \(A=132\) is spatially more extended.
    \item Wave functions become slightly broader.
    \item Overlap integrals are modified.
    \item Two-body matrix elements change through the basis parameter.
\end{itemize}
\end{frame}

%------------------------------------------------
\begin{frame}[t]{Phenomenological Scaling of Matrix Elements}
\small
A commonly used phenomenological scaling law for shell-model two-body
matrix elements is
\[
V_{JT}(A)=V_{JT}(A_0)\left(\frac{A}{A_0}\right)^{-\alpha},
\qquad \alpha\approx 0.3 .
\]

For the present example:
\[
V(A=132)=V(A=100)\left(\frac{132}{100}\right)^{-0.3} .
\]

This gives
\[
V(A=132)\approx 0.92\,V(A=100) .
\]
\end{frame}

%------------------------------------------------
\begin{frame}[t]{Physical Interpretation}
\small
Why does the interaction strength decrease slightly?

\medskip

As the nucleus becomes larger:
\begin{itemize}
    \item the oscillator length \(b\) increases,
    \item single-particle wave functions spread out,
    \item overlap integrals decrease slightly,
    \item effective interaction matrix elements become weaker.
\end{itemize}

\medskip

The scaling approximately captures changes in nuclear radius, basis
dependence, renormalization effects, and effective many-body correlations.
\end{frame}

%------------------------------------------------
\begin{frame}[t]{More Fundamental Treatment}
\small
The phenomenological scaling law is only approximate.

\medskip

A more rigorous approach would:
\begin{enumerate}
    \item recompute the harmonic oscillator basis with the new \(b\) value,
    \item recalculate Talmi integrals,
    \item reevaluate the two-body matrix elements.
\end{enumerate}

\medskip

This is especially important for finite-range interactions, chiral EFT
interactions, large changes in mass number, and precision spectroscopy.
\end{frame}

%------------------------------------------------
\begin{frame}[t]{Talmi--Moshinsky Formalism}
\small
Two-body matrix elements in the oscillator basis are often expressed as
\[
\langle ab|V|cd\rangle = \sum_N C_N I_N .
\]

Here
\begin{itemize}
    \item \(C_N\) are Moshinsky brackets,
    \item \(I_N\) are Talmi integrals.
\end{itemize}

The Talmi integrals depend explicitly on the oscillator length:
\[
I_N=I_N(b) .
\]

Therefore changing \(A\) changes the matrix elements through the basis
parameter \(b\).
\end{frame}

%------------------------------------------------
\begin{frame}[t]{Example: Scaling from \(A=100\) to \(A=132\)}
\small
Using the phenomenological prescription
\[
V(A)=V(A_0)\left(\frac{A}{A_0}\right)^{-0.3},
\]
we obtain
\[
V(132)=V(100)\left(\frac{132}{100}\right)^{-0.3} .
\]

Numerically,
\[
\left(\frac{132}{100}\right)^{-0.3}\approx 0.92 .
\]

Hence
\[
V(132)\approx 0.92\,V(100),
\]
which corresponds to a reduction of approximately \(8\%\).
\end{frame}

%------------------------------------------------
\begin{frame}[t]{Typical Usage in Shell-Model Calculations}
\small
Mass scaling is commonly employed in:
\begin{itemize}
    \item USD interactions,
    \item large-scale shell-model calculations,
    \item effective interactions derived from \(G\)-matrix theory.
\end{itemize}

\medskip

Advantages:
\begin{itemize}
    \item simple implementation,
    \item captures average trends,
    \item computationally inexpensive.
\end{itemize}
\end{frame}

%------------------------------------------------
\begin{frame}[t]{Limitations of Simple Mass Scaling}
\small
The same scaling has important limitations:
\begin{itemize}
    \item it is not fully microscopic,
    \item it cannot replace full renormalization,
    \item it may fail far from fitted mass regions,
    \item it cannot capture detailed monopole and multipole retuning.
\end{itemize}

\medskip

For high-precision spectroscopy, the preferred procedure is to recompute or
refit the effective interaction in the relevant mass region and oscillator
basis.
\end{frame}

%------------------------------------------------
\begin{frame}[t]{Summary: Scaling Relations}
\small
The central scaling relations are:
\[
\hbar\omega \propto A^{-1/3},
\qquad
b\propto A^{1/6} .
\]

Effective shell-model matrix elements are often scaled phenomenologically as
\[
V(A)=V(A_0)\left(\frac{A}{A_0}\right)^{-0.3} .
\]

For scaling from \(A=100\) to \(A=132\), this gives
\[
V(132)\approx 0.92\,V(100) .
\]
\end{frame}

%------------------------------------------------
\begin{frame}[t]{Summary: Practical Recommendation}
\small
A simple first estimate is therefore to multiply the \(A=100\) two-body
matrix elements by approximately \(0.92\).

\medskip

However, a more accurate treatment requires recomputing the matrix elements
using the new oscillator basis. Reference~\cite{MHJ1995} shows how this can
be done using many-body theories.
\end{frame}

%------------------------------------------------
\begin{frame}[t]{References}
\footnotesize
\begin{thebibliography}{99}

\bibitem{Talmi}
I. Talmi,
\newblock \emph{Simple Models of Complex Nuclei},
Harwood Academic Publishers (1993).

\bibitem{MHJ1995}
M. Hjorth-Jensen, T.~T.~S. Kuo, and E. Osnes,
\newblock \emph{Realistic effective interactions for nuclear systems},
Physics Reports \textbf{261}, 125 (1995).

\end{thebibliography}
\end{frame}

%------------------------------------------------
\end{document}
